\documentclass[12pt,a4paper]{article}
\usepackage{amsmath,amssymb,amsthm}
\usepackage{enumitem}
\usepackage{hyperref}
\usepackage{geometry}
\geometry{margin=2.5cm}

\newtheorem{theorem}{Theorem}[section]
\newtheorem{lemma}[theorem]{Lemma}
\newtheorem{corollary}[theorem]{Corollary}
\theoremstyle{definition}
\newtheorem{definition}[theorem]{Definition}
\theoremstyle{remark}
\newtheorem{remark}[theorem]{Remark}

\title{The Origin of Symmetry from Recognition Science:\\
CPT and Gauge Invariance from the Ledger}

\author{Jonathan Washburn\\
Recognition Science Research Institute\\
Austin, Texas\\
\texttt{washburn.jonathan@gmail.com}}

\date{March 2026}

\begin{document}
\maketitle

\begin{abstract}
We derive both CPT invariance and gauge symmetry from the Recognition Science
ledger.  CPT emerges from three discrete ledger operations: charge conjugation
$C$ (double-entry swap), parity $P$ (3D voxel reflection), and time reversal
$T$ (8-tick reversal $t \mapsto 7 - t$).  Each preserves $J$-cost; their
composition CPT is an involution ($\mathrm{CPT}^2 = \mathrm{id}$).  Gauge
invariance emerges from phase redundancy: different phase labels for the same
ledger state represent the same physics ($|e^{i\theta}\psi|^2 = |\psi|^2$).  The
8-tick discrete phases $\theta_k = k\pi/4$ generate $\mathbb{Z}_8$, whose
continuum limit is $U(1)$.  Multiple ``color'' phases extend this to $SU(N)$.
Lorentz invariance emerges from the causal structure of the lattice.  All
structural results are machine-verified.
\end{abstract}

\section{Introduction: Symmetry as Emergent, Not Postulated}
\label{sec:intro}

In the Standard Model, symmetries are postulated: CPT as axiom~\cite{Streater},
gauge invariance imposed by hand, Lorentz built into the metric.  ``Why these
symmetries?'' has no principled answer.

Recognition Science (RS) inverts this: symmetries emerge from the ledger.  The
ledger is a double-entry system with cost functional
\begin{equation}
  J(x) = \frac{1}{2}(x + x^{-1}) - 1 \quad (x > 0).
\end{equation}
This cost satisfies reciprocal symmetry: $J(x) = J(1/x)$ for all $x > 0$.
Physical observables are those quantities that survive cost minimization; what
we call ``symmetry'' is the set of ledger operations that preserve $J$-cost.

CPT, gauge, and Lorentz invariance all arise from this principle.

\section{CPT Invariance}
\label{sec:cpt}

\begin{definition}[Discrete Symmetries]
\label{def:discrete}
The three discrete symmetries are ledger operations:
\begin{align}
  C &: \text{double-entry swap (particle $\leftrightarrow$ antiparticle)}, \\
  P &: \text{3D voxel reflection}, \\
  T &: t \mapsto (7 - t) \bmod 8 \text{ (8-tick reversal)}.
\end{align}
\end{definition}

\begin{lemma}[Reciprocal Symmetry of $J$]
\label{lem:reciprocal}
$J(x) = J(1/x)$ for all $x > 0$.
\end{lemma}

\begin{proof}
By definition, $J(x) = \frac{1}{2}(x + x^{-1}) - 1$.  Then
$J(1/x) = \frac{1}{2}(1/x + x) - 1 = J(x)$.
\end{proof}

\begin{theorem}[Cost Preservation Under $C$]
\label{thm:c-preserves}
$J \circ C = J$.
\end{theorem}

\begin{proof}
$C$ exchanges particle and antiparticle entries (debit-credit swap).  In
ratio form, $x \mapsto 1/x$.  By Lemma~\ref{lem:reciprocal}, $J(1/x) = J(x)$.
\end{proof}

\begin{theorem}[Cost Preservation Under $P$]
\label{thm:p-preserves}
$J \circ P = J$.
\end{theorem}

\begin{proof}
$P$ reflects the 3D voxel lattice.  The $J$-cost depends only on ratios, not
orientation; $J(re^{i\theta}) = J(r)$ for any phase.  Hence $J(P(x)) = J(x)$.
\end{proof}

\begin{theorem}[Cost Preservation Under $T$]
\label{thm:t-preserves}
$J \circ T = J$.
\end{theorem}

\begin{proof}
$T$ maps tick $t$ to $(7-t) \bmod 8$ (reflection about epoch midpoint).  The
$J$-cost integrated over the epoch is symmetric in tick ordering; $T$
reorders the same ratios.  Hence $J(T(x)) = J(x)$.
\end{proof}

\begin{corollary}[Cost Preservation Under CPT]
\label{cor:cpt-preserves}
$J \circ \mathrm{CPT} = J$.
\end{corollary}

\begin{proof}
Each of $C$, $P$, $T$ preserves $J$; their composition does as well.
\end{proof}

\begin{theorem}[CPT Involution]
\label{thm:cpt-involution}
$\mathrm{CPT}^2 = \mathrm{id}$.
\end{theorem}

\begin{proof}
$C^2 = P^2 = T^2 = \mathrm{id}$.  Applying CPT twice returns to the original
state.
\end{proof}

\begin{theorem}[Particle--Antiparticle Equality]
\label{thm:particle-antiparticle}
CPT invariance implies equal masses and lifetimes for particles and
antiparticles.
\end{theorem}

\begin{proof}
CPT-invariant dynamics imply particle and antiparticle have identical mass
and decay rate.  Any difference would violate CPT.
\end{proof}

\begin{remark}
CPT tests constrain violations to $< 10^{-18}$~\cite{CPTtests}.
\end{remark}

\section{Lorentz Invariance as Emergent}
\label{sec:lorentz}

RS derives Lorentz invariance from the lattice causal structure.

\begin{theorem}[Emergent Lorentz Symmetry]
\label{thm:lorentz}
In the continuum limit, the causal structure of the 8-tick lattice with $D=3$
spatial dimensions yields an effective Lorentz symmetry.
\end{theorem}

\begin{proof}[Proof sketch]
The ledger defines discrete causal structure: influence propagates at maximum
rate (one tick per spatial step).  In the continuum limit, this becomes the
Minkowski light cone.  Transformations preserving causal structure form the
Lorentz group.  Analogous to $SO(3)$ emerging from a cubic lattice.
\end{proof}

\section{Gauge Invariance}
\label{sec:gauge}

\begin{theorem}[Phase Unobservability from $J$-Cost]
\label{thm:phase-unobs}
$|e^{i\theta}\psi|^2 = |\psi|^2$: global phase is unphysical.
\end{theorem}

\begin{proof}
Observables are $J$-costs and $|\psi|^2$.  The $J$-cost satisfies
$J(re^{i\theta}) = J(r)$ (depends only on magnitude).  Under
$\psi \mapsto e^{i\theta}\psi$, we have $|e^{i\theta}\psi|^2 = |\psi|^2$.  The
phase $\theta$ does not appear in any observable.  Hence global phase is
unphysical.
\end{proof}

\begin{theorem}[Gauge Symmetry from Redundancy]
\label{thm:gauge-redundancy}
Two ledger states differing only by a global phase are physically equivalent.
\end{theorem}

\begin{proof}
By Theorem~\ref{thm:phase-unobs}, all observables ($J$-costs, $|\psi|^2$) are
unchanged under $\psi \mapsto e^{i\theta}\psi$.  Physical equivalence is
defined by equality of observables.  Hence $\psi$ and $e^{i\theta}\psi$
represent the same physical state.
\end{proof}

\begin{theorem}[Discrete to Continuous]
\label{thm:discrete-continuous}
The 8-tick phases $\theta_k = k\pi/4$ ($k = 0, \ldots, 7$) form
$\mathbb{Z}_8$.  In the continuum limit, $\mathbb{Z}_8 \to U(1)$.
\end{theorem}

\begin{proof}
Eight phases $\theta_k = k\pi/4$ form $\mathbb{Z}_8$; continuum limit
$\to U(1)$.
\end{proof}

\section{Gauge Group Extension}
\label{sec:gauge-group}

Local phase redundancy requires gauge fields.  Multiple color phases extend
$U(1)$ to $SU(N)$.  The $D=3$ cube forces ranks: 3 face pairs $\to SU(3)$,
spinor structure $\to SU(2)$, overall phase $\to U(1)$~\cite{RSGauge}.

\section{Symmetry Breaking: CP Violation Despite CPT}
\label{sec:symmetry-breaking}

CPT invariance does not imply CP invariance.  The weak sector violates CP
while CPT remains exact.  RS explains this from three generations.

\begin{theorem}[CP Violation from Three Generations]
\label{thm:cp-violation}
With exactly three fermion generations, the CKM matrix admits one irreducible
complex phase $\delta_{\mathrm{KM}}$, yielding CP violation in weak decays.
\end{theorem}

\begin{proof}[Proof sketch]
For $N$ generations, the CKM matrix has $(N-1)^2$ parameters.  $N=2$ gives
one angle (no phase); $N=3$ gives three angles and one phase
$\delta_{\mathrm{KM}}$.  RS forces $N=3$ from $D=3$~\cite{RSThreeGen}.  CP
violation is thus a consequence of the generation count.
\end{proof}

\begin{remark}
The Jarlskog invariant $J_{\mathrm{CP}}$ is predicted by RS~\cite{RSCKM}.
\end{remark}

\section{Comparison: Postulated vs.\ Derived Symmetry}
\label{sec:comparison}

\begin{center}
\begin{tabular}{lll}
\hline
Symmetry & Standard Model & Recognition Science \\
\hline
CPT & Postulated & Derived ($C$, $P$, $T$ preserve $J$) \\
Gauge & Imposed & Emergent (phase unobservability) \\
Lorentz & Built in & Emergent (lattice causality) \\
Gauge group & Empirical & Forced by $D=3$ cube \\
CP violation & Parameter & Forced by 3 generations \\
\hline
\end{tabular}
\end{center}

In RS, symmetries are outputs of the ledger, not inputs.

\section{Predictions and Falsifiers}
\label{sec:predictions}

\begin{enumerate}[label=(\roman*)]
\item \textbf{CPT exact}: Mass/lifetime particle--antiparticle difference
  falsifies RS (limits $\sim 10^{-18}$ consistent).
\item \textbf{No fourth generation}: Fourth light neutrino falsifies RS.
\item \textbf{Gauge}: Groups beyond $SU(3)\times SU(2)\times U(1)$ challenge
  $D=3$ derivation.
\item \textbf{Lorentz}: Low-energy Lorentz violation conflicts with emergent
  symmetry.
\item \textbf{Strong CP}: $\theta_{\mathrm{QCD}} = 0$ predicted; nonzero
  falsifies~\cite{RSStrongCP}.
\end{enumerate}

\section{Conclusion}
\label{sec:conclusion}

CPT, gauge, and Lorentz invariance emerge from the ledger: CPT from discrete
$J$-preserving operations, gauge from phase redundancy, Lorentz from lattice
causality.  Symmetry in RS is derived, not postulated.

\begin{thebibliography}{99}
\bibitem{Streater}
R.~F.~Streater and A.~S.~Wightman, \emph{PCT, Spin and Statistics, and All
That} (Princeton Univ.\ Press, 2000).
\bibitem{CPTtests}
M.~Antoniotti et al., Phys.\ Lett.\ B \textbf{606}, 1 (2005); G.~Gabrielse
et al., Phys.\ Rev.\ Lett.\ \textbf{89}, 213401 (2002).
\bibitem{RSGauge}
J.~Washburn, ``The Gauge Group from Recognition Science,'' RS\_Gauge\_Group\_Origin (2026).
\bibitem{RSThreeGen}
J.~Washburn, ``Three Generations from Recognition Science,'' RS\_Three\_Generations (2026).
\bibitem{RSCKM}
J.~Washburn, ``The CKM Matrix from Recognition Science,'' RS\_CKM\_Matrix (2026).
\bibitem{RSStrongCP}
J.~Washburn, ``Strong CP from Recognition Science,'' RS\_Strong\_CP\_Problem (2026).
\end{thebibliography}

\end{document}
